% GENERATED FILE. Edit canonical lessons, then run npm run generate:books.
% canonical-source-hash: fnv1a64:b1e0e11c723a057b
% canonical-lessons: ES-C471-demasiado, ES-C471-nivel, ES-C471-avanzado, ES-C471-principiante, ES-C471-sencillo, ES-C471-perderse, ES-R471-repaso-curso, ES-C471-sintesis-curso

\chapter{El curso de cocina}
\label{ch:es-curso}
\hlchaptermodality{471}

\begin{chapteropening}
\textbf{By the end of this chapter:} \emph{I can follow somebody explaining that a class did not suit them: hear that the level was too high, that they kept losing the thread, and that they moved to a beginners' group, and say the same about something of my own without criticising anyone.}

Follow two friends on a cooking class that was too advanced, then say the same about something of your own.
\end{chapteropening}

\section[demasiado]{demasiado — too much}

\label{lesson:ES-C471-demasiado}



\begin{quote}
Say \textbf{bastante} and \textbf{poco}. You have a scale with a bottom and a middle. Today it gets its top.
\end{quote}

\subsection*{You'll want to know: demasiado}
\textbf{demasiado} — too much, too many, too.

\begin{quote}
El curso era \textbf{demasiado} difícil para mí.
\end{quote}

Before an adjective it does not change: \emph{demasiado difícil}, \emph{demasiado avanzada}. Before a noun it agrees like any adjective: \emph{demasiada gente}, \emph{demasiados libros}.

\begin{grammarlens}[title={Grammar Lens: the word that says it went past the point}]
Bastante built you a scale of degree. This is the member above the top of it, and the step matters.

\noindent
\begin{tabularx}{\linewidth}{@{}>{\raggedright\arraybackslash}X>{\raggedright\arraybackslash}X@{}}
\toprule
\textbf{the word} & \textbf{where on the scale} \\
\midrule
poco & below what is wanted \\
bastante & enough — the point itself \\
\textbf{demasiado} & past the point \\
\bottomrule
\end{tabularx}

\emph{Muy difícil} is a judgement about the thing. \textbf{Demasiado difícil} is a judgement about a \textbf{fit}: too hard \emph{for something}, usually for you. That is why it so often arrives with \emph{para} — \emph{demasiado avanzado para mí}.

English blurs this, because \emph{too} and \emph{very} sit close together in speech. Spanish does not: say \emph{muy} when you mean intense, and \emph{demasiado} when you mean it has gone past what works.
\end{grammarlens}

\begin{cousinweb}
Underneath is \textbf{de más} — \emph{of more}, in excess — which became \emph{demasía}, an excess, and then this adjective. The \emph{más} in it is the \emph{más} you already use, from Latin \textbf{magis}, more.

So the word is built out of an admission: there is \textbf{more} here than the situation has room for. Not more than exists, more than \textbf{fits}.

Keep that and the grammar follows. \emph{Demasiado} always implies a measure it has overrun, even when nobody says what the measure was.
\end{cousinweb}

\subsection*{Your turn}
Say these aloud:

\begin{itemize}
\raggedright
  \item it was too difficult for me
  \item the difference between \emph{muy difícil} and \emph{demasiado difícil}
  \item too many books, with the agreement
\end{itemize}

\begin{tcolorbox}[breakable,colback=teal!4,colframe=teal!35!black,title={Before you move on}]
Too much? (\textbf{Demasiado}.) Which taught word is inside it? (\emph{\textbf{Más}}.) Does it mean \emph{very}? (\textbf{No — past the point that fits}.)
\end{tcolorbox}
\section[el nivel]{el nivel — the level}

\label{lesson:ES-C471-nivel}



\begin{quote}
Say \textbf{el curso} and \textbf{demasiado}. A class has one of these, and when it is wrong the second word is how you say so.
\end{quote}

\subsection*{You'll want to know: el nivel}
\textbf{el nivel} — the level.

\begin{quote}
El \textbf{nivel} era demasiado alto para ella.
\end{quote}

\emph{Nivel} is what a class, a language exam or a game has. \emph{¿Qué nivel tienes?} is the ordinary way to ask how far along somebody is. Note the stress: \emph{ni-VEL}, on the end, and the plural is \textbf{los niveles}.

It is also a level in the physical sense — a floor of a building, the surface of water, the tool a builder uses.

\begin{grammarlens}[title={Grammar Lens: the phrase the question will use}]
Two ways of saying the same complaint, and you need to recognise both.

\noindent
\begin{tabularx}{\linewidth}{@{}>{\raggedright\arraybackslash}X>{\raggedright\arraybackslash}X@{}}
\toprule
\textbf{about the class} & \textbf{about its level} \\
\midrule
el curso era \textbf{demasiado avanzado} & el \textbf{nivel} era \textbf{demasiado alto} \\
\bottomrule
\end{tabularx}

A speaker will usually say the first. Somebody reporting it afterwards — a summary, a question, an exam option — will often reach for the second, because \emph{nivel} lets you talk about the difficulty without naming the class again.

\emph{Alto} and \emph{bajo} are what \emph{nivel} takes, not \emph{grande} or \emph{mucho}. A level is a \textbf{height}.
\end{grammarlens}

\begin{cousinweb}
El litro already told you about Latin \textbf{libra}, the balance — hiding in \emph{lb}, in \textbf{Libra} the scales in the sky, in \textbf{deliberate} and \textbf{equilibrium}.

\emph{Nivel} is the same word, made small. Latin \textbf{libella} is a \emph{little libra}, a builder's level: a tiny balance that tells you when a surface is true. It reached Spanish through French, and on the way the first \textbf{l} turned into an \textbf{n} — \emph{livel} became \emph{nivel} — because two \emph{l} sounds so close together pull apart.

English took \emph{libella} too, without that swap, and got \textbf{level}.

So \emph{nivel} and \emph{litro} are cousins, and a level is not a measure of height so much as an instrument for finding \textbf{balance}.
\end{cousinweb}

\subsection*{Your turn}
Say these aloud:

\begin{itemize}
\raggedright
  \item the level was too high for her
  \item ask somebody what level they have
  \item which taught word \emph{nivel} is a cousin of
\end{itemize}

\begin{tcolorbox}[breakable,colback=teal!4,colframe=teal!35!black,title={Before you move on}]
The level? (\textbf{El nivel}.) Plural? (\textbf{Los niveles}.) What Latin word is underneath? (\emph{\textbf{Libra}}, the balance — \emph{libella}, a little one.)
\end{tcolorbox}
\section[avanzado]{avanzado — advanced}

\label{lesson:ES-C471-avanzado}



\begin{quote}
Say \textbf{el curso} and \textbf{antes}. The second is hiding inside today's word, and you would not guess it from the outside.
\end{quote}

\subsection*{You'll want to know: avanzado}
\textbf{avanzado} — advanced.

\begin{quote}
El curso de los lunes era demasiado \textbf{avanzado} para mí.
\end{quote}

\emph{Un nivel avanzado}, \emph{un grupo avanzado}: it describes the class, not the person. For the person you say \emph{tiene un nivel avanzado}.

It agrees like any adjective — \emph{avanzada}, \emph{avanzados}, \emph{avanzadas}.

\begin{grammarlens}[title={Grammar Lens: a participle doing an adjective's work}]
\emph{Avanzado} is the past participle of \emph{avanzar}, to advance, standing in as an adjective. Spanish does this constantly, and you have seen it before:

\noindent
\begin{tabularx}{\linewidth}{@{}>{\raggedright\arraybackslash}X>{\raggedright\arraybackslash}X@{}}
\toprule
\textbf{the participle} & \textbf{what it describes} \\
\midrule
\textbf{invitado} & the one who was invited \\
\textbf{avanzado} & the one that has moved ahead \\
\bottomrule
\end{tabularx}

El invitado made the point first — Spanish names a person by what was done to them. Here the same move names a \textbf{course} by how far it has got.

So a \emph{curso avanzado} is not a course that is difficult in itself. It is one that has already travelled, and expects you to have travelled with it. That distinction is the whole reason the woman in the recording changed groups.
\end{grammarlens}

\begin{cousinweb}
\textbf{Avanzar} is Latin \textbf{ab ante} — \emph{from in front} — which became \emph{abantiare} and then the Spanish verb. The \emph{ante} in it is the same \emph{ante} you already hold in \textbf{antes}.

To advance is to put yourself \textbf{in front}. Not faster; further along.

One thing worth knowing about the English. \textbf{Advance} has a \textbf{d} that does not belong to it: the word arrived as \emph{avancer}, and English later added the \emph{d} on the mistaken assumption that the \emph{av-} was a worn-down Latin \emph{ad-}. It was not. Spanish never made that mistake, which is why \emph{avanzar} has no \emph{d} and never did.
\end{cousinweb}

\subsection*{Your turn}
Say these aloud:

\begin{itemize}
\raggedright
  \item the Monday course was too advanced for me
  \item which taught word is hiding inside \emph{avanzar}
  \item why English \emph{advance} has a d and Spanish has none
\end{itemize}

\begin{tcolorbox}[breakable,colback=teal!4,colframe=teal!35!black,title={Before you move on}]
Advanced? (\textbf{Avanzado}.) What is it the participle of? (\textbf{Avanzar}.) What two Latin words are inside it? (\emph{\textbf{Ab ante}} — from in front.)
\end{tcolorbox}
\section[el principiante]{el principiante — the beginner}

\label{lesson:ES-C471-principiante}



\begin{quote}
Say \textbf{el principio}. Say \textbf{el estudiante}. One gives today's word its meaning and the other gives it its ending — try it before reading on.
\end{quote}

\subsection*{You'll want to know: el principiante}
\textbf{el principiante} — the beginner.

\begin{quote}
Ahora voy con los \textbf{principiantes} y disfruto mucho más.
\end{quote}

\textbf{La principiante} for a woman. It is the natural opposite of what you learned a moment ago:

\noindent
\begin{tabularx}{\linewidth}{@{}>{\raggedright\arraybackslash}X>{\raggedright\arraybackslash}X@{}}
\toprule
\textbf{the class} & \textbf{the person in it} \\
\midrule
un curso \textbf{avanzado} & un \textbf{principiante} \\
\bottomrule
\end{tabularx}

\begin{grammarlens}[title={Grammar Lens: you could have built this one}]
El estudiante gave you the rule outright: \textbf{verb plus -ante names the person doing it.} \emph{Estudiar}, cut the \emph{-ar}, add \emph{-ante}.

Here it is again, one step further back. From \textbf{el principio}, the beginning, Spanish makes the verb \textbf{principiar}, to begin. Take that, cut the \emph{-ar}, add \emph{-ante}:

\begin{quote}
princip\textbf{iar} $\to$ princip\textbf{iante}
\end{quote}

The one who is beginning. The ending has not changed its job since chapter nine; the only new thing is the verb it is sitting on.

English wore the same ending down to \emph{-ant}, which is why \emph{participant}, \emph{assistant} and \emph{inhabitant} all name the doer too.
\end{grammarlens}

\begin{cousinweb}
\emph{El principio} already traced this for you: Latin \textbf{principium}, from \textbf{princeps}, the one who takes first place — the same first-taking inside English \textbf{prince}, \textbf{principal} and \textbf{principle}.

So a \emph{principiante} is not somebody who is bad at the thing. The word says only that they are at the \textbf{start} of it. Spanish has kept that neutral, and so should you when you use it about yourself.
\end{cousinweb}

\subsection*{Your turn}
Say these aloud:

\begin{itemize}
\raggedright
  \item now I go with the beginners
  \item build the word from \emph{principiar} out loud
  \item the opposite of \emph{un curso avanzado}
\end{itemize}

\begin{tcolorbox}[breakable,colback=teal!4,colframe=teal!35!black,title={Before you move on}]
The beginner? (\textbf{El principiante}.) What does \emph{-ante} add? (\textbf{The one who does it}.) Does the word say they are bad at it? (\textbf{No — only that they are at the start}.)
\end{tcolorbox}
\section[sencillo]{sencillo — simple}

\label{lesson:ES-C471-sencillo}



\begin{quote}
Say \textbf{difícil}. Say \textbf{el principiante}. Today's word is what a beginner's class is full of, and it is not the opposite of the first one.
\end{quote}

\subsection*{You'll want to know: sencillo}
\textbf{sencillo} — simple, plain, straightforward.

\begin{quote}
Hacemos cosas \textbf{sencillas}, pero las hacemos bien.
\end{quote}

\emph{Una comida sencilla} — a plain meal. \emph{Un billete sencillo} — a one-way ticket, where the word means \emph{single} rather than \emph{easy}.

\begin{grammarlens}[title={Grammar Lens: sencillo is not fácil}]
English lets \emph{simple} and \emph{easy} stand in for each other. Spanish keeps them apart, and the line runs between \textbf{parts} and \textbf{effort}.

\noindent
\begin{tabularx}{\linewidth}{@{}>{\raggedright\arraybackslash}X>{\raggedright\arraybackslash}X@{}}
\toprule
\textbf{the word} & \textbf{what it measures} \\
\midrule
\textbf{sencillo} & how few parts it has \\
fácil & how little effort it costs \\
\bottomrule
\end{tabularx}

A dish with three ingredients is \emph{sencillo}. A dish you can cook half asleep is \emph{fácil}. They often overlap and they are not the same claim — climbing a staircase is \emph{sencillo} and not remotely \emph{fácil}.

So \emph{hacemos cosas sencillas, pero las hacemos bien} is not modest about the effort. It says the things have few parts, and then insists the execution is good. Hear it as \emph{easy} and you lose the point the speaker is making.
\end{grammarlens}

\begin{cousinweb}
Latin \textbf{singulus}, one at a time, single — through a late form \emph{singellus}. English took the same root for \textbf{single}, \textbf{singular} and \textbf{singly}.

That is why \emph{un billete sencillo} means a \textbf{single}: the word's oldest sense is \emph{one}, not \emph{easy}. Something \emph{sencillo} is made of one thing, or nearly.

The \emph{-illo} on the end is not doing what it usually does. Elsewhere it makes things small — but here it was already fused into the Latin, and there is no word \emph{senci} underneath for it to shrink.
\end{cousinweb}

\subsection*{Your turn}
Say these aloud:

\begin{itemize}
\raggedright
  \item we do simple things, but we do them well
  \item the difference between \emph{sencillo} and \emph{fácil}
  \item what \emph{un billete sencillo} is
\end{itemize}

\begin{tcolorbox}[breakable,colback=teal!4,colframe=teal!35!black,title={Before you move on}]
Simple? (\textbf{Sencillo}.) Does it mean easy? (\textbf{No — few parts, not little effort}.) What Latin word is underneath? (\emph{\textbf{Singulus}} — one at a time.)
\end{tcolorbox}
\section[perderse]{perderse — to get lost}

\label{lesson:ES-C471-perderse}



\begin{quote}
Say \textbf{perder}. Say \textbf{equivocarse}. You have had both, and the second one already told you what happens when you put \emph{-se} on the first.
\end{quote}

\subsection*{You'll want to know: perderse}
\textbf{perderse} — to get lost.

\begin{quote}
Todos habían hecho cursos antes y yo \textbf{me perdía}.
\end{quote}

Note what is lost there. Not the way — the \textbf{thread}. In a class, in a film, in a conversation, \emph{me pierdo} is \emph{I lose track}, and it is the most useful sense the verb has for you.

\emph{Me he perdido} — I've got lost, or I've lost the thread. \emph{No te pierdas} — don't get lost, and also don't miss it.

\begin{grammarlens}[title={Grammar Lens: the family you were told to notice}]
Equivocarse named this verb explicitly, in a group worth recalling whole:

\noindent
\begin{tabularx}{\linewidth}{@{}>{\raggedright\arraybackslash}X>{\raggedright\arraybackslash}X@{}}
\toprule
\textbf{plain} & \textbf{reflexive} \\
\midrule
perder — to lose a thing & \textbf{perderse} — to get lost \\
olvidar — to forget a thing & olvidarse — it slips your mind \\
\bottomrule
\end{tabularx}

The reflexive marks something that \textbf{befalls you} rather than something you do to an object. \emph{Perder el paraguas} is an act with a victim. \emph{Perderse} has no object at all: the losing happens to you.

That is why \emph{me perdía} is not an admission of carelessness. It reports a state she was in, which is exactly why it is the honest thing to say about a class that had moved on without her.
\end{grammarlens}

\begin{cousinweb}
Latin \textbf{perdere}: \textbf{per-} plus \textbf{dare}, to give — to give \textbf{away} thoroughly, to hand something over for good.

That \emph{per-} is the thorough one, the same intensifier inside \textbf{perfect} (thoroughly made) and \textbf{perturb}. So losing, in Latin, is a kind of extreme giving: the thing goes and does not come back.

English \textbf{perdition} is the same word, kept for the most complete loss there is.
\end{cousinweb}

\subsection*{Your turn}
Say these aloud:

\begin{itemize}
\raggedright
  \item everyone had done courses before and I kept losing the thread
  \item the difference between \emph{perder} and \emph{perderse}
  \item what \emph{per-} adds to \emph{dare}
\end{itemize}

\begin{tcolorbox}[breakable,colback=teal!4,colframe=teal!35!black,title={Before you move on}]
To get lost? (\textbf{Perderse}.) In a class, what do you lose? (\textbf{The thread}.) Why the \emph{-se}? (\textbf{It happens to you, not to an object}.)
\end{tcolorbox}
\section[(demasiado, el nivel, avanzado, el …]{(demasiado, el nivel, avanzado, el principiante, sencillo, perderse) — from cold}

\label{lesson:ES-R471-repaso-curso}



\begin{quote}
Six words, no prompts. How far past the point it was, what it was too high a one of, what the hard class was, who is in the other one, what they cook there, and what kept happening to you in the first.
\end{quote}

\subsection*{You'll want to know: the six, cold}
\noindent
\begin{tabularx}{\linewidth}{@{}>{\raggedright\arraybackslash}X>{\raggedright\arraybackslash}X@{}}
\toprule
\textbf{English} & \textbf{Spanish} \\
\midrule
too much & \textbf{demasiado} \\
the level & \textbf{el nivel} \\
advanced & \textbf{avanzado} \\
the beginner & \textbf{el principiante} \\
simple & \textbf{sencillo} \\
to get lost & \textbf{perderse} \\
\bottomrule
\end{tabularx}

\begin{grammarlens}[title={Grammar Lens: three endings, all of them already yours}]
Not one ending in this chapter is new. Each was taught somewhere else and is here doing its ordinary job.

\noindent
\begin{tabularx}{\linewidth}{@{}>{\raggedright\arraybackslash}X>{\raggedright\arraybackslash}X@{}}
\toprule
\textbf{the ending} & \textbf{what it does} \\
\midrule
\textbf{-ado} & a participle used as an adjective \\
\textbf{-ante} & the one who does it \\
\textbf{-se} & it happens to you \\
\bottomrule
\end{tabularx}

\emph{El invitado} gave the first, \emph{el estudiante} the second, \emph{equivocarse} the third. So three of these six words were readable before you were told, and the chapter's real work was elsewhere: knowing which \textbf{sense} each carries, and filling two holes the scale had left.
\end{grammarlens}

\begin{cousinweb}
Four of the six hide a word you already hold.

\noindent
\begin{tabularx}{\linewidth}{@{}>{\raggedright\arraybackslash}X>{\raggedright\arraybackslash}X@{}}
\toprule
\textbf{word} & \textbf{underneath} \\
\midrule
demasiado & \emph{de más} — the \textbf{más} you know \\
avanzado & \emph{ab ante} — the \textbf{antes} you know \\
principiante & \emph{principium} — the \textbf{principio} you know \\
\bottomrule
\end{tabularx}

And \textbf{el nivel} is Latin \emph{libella}, a little \emph{libra} — so it is a cousin of \textbf{el litro}, which carried that balance for you already.

\emph{Sencillo} is Latin \emph{singulus}, one at a time — English \textbf{single}. And \emph{perderse} is \emph{per-} plus \emph{dare}, to give away thoroughly, which English kept in \textbf{perdition}.
\end{cousinweb}

\subsection*{Your turn}
Say these aloud:

\begin{itemize}
\raggedright
  \item the Monday class was too advanced and I kept losing the thread
  \item now I go with the beginners and we do simple things
  \item which of the six means few parts rather than little effort
\end{itemize}

\begin{tcolorbox}[breakable,colback=teal!4,colframe=teal!35!black,title={Before you move on}]
Too much? (\textbf{Demasiado}.) The level? (\textbf{El nivel}.) Advanced? (\textbf{Avanzado}.) The beginner? (\textbf{El principiante}.) Simple? (\textbf{Sencillo}.) To get lost? (\textbf{Perderse}.)
\end{tcolorbox}
\section[Practice]{(el curso de cocina) — follow it, then say your own}

\label{lesson:ES-C471-sintesis-curso}



\begin{quote}
Say \textbf{avanzado} and \textbf{perderse}. The first explains the second in everything below.
\end{quote}

\begin{tcolorbox}[breakable,skin=enhanced,colback=orange!6,colframe=orange!45!black,boxrule=0.5pt,arc=1mm,left=6pt,right=6pt,top=4pt,bottom=4pt,fonttitle=\bfseries,title={Read this}]
\begin{quote}
— ¿Qué tal el curso de cocina? ¿Sigues yendo? — Sí, pero he cambiado de grupo. El de los lunes era demasiado \textbf{avanzado}   para mí; todos habían hecho cursos antes y yo me \textbf{perdía}. — ¿Y ahora? — Ahora voy los miércoles con los \textbf{principiantes} y disfruto mucho más.   Hacemos cosas \textbf{sencillas}, pero las hacemos bien.
\end{quote}

\textbf{El de los lunes.} The noun \emph{grupo} is dropped and \emph{el} carries it — the same move the printer notice made with \emph{la de la tercera planta}.
\end{tcolorbox}

\begin{grammarlens}[title={Grammar Lens: she never says the class was bad}]
Every word of the complaint is about \textbf{fit}, not quality.

\noindent
\begin{tabularx}{\linewidth}{@{}>{\raggedright\arraybackslash}X>{\raggedright\arraybackslash}X@{}}
\toprule
\textbf{what she says} & \textbf{what she does not say} \\
\midrule
\emph{demasiado avanzado para mí} & that the class was bad \\
\emph{yo me perdía} & that the teacher was bad \\
\bottomrule
\end{tabularx}

\emph{Para mí} puts the mismatch between the class and her; \emph{me perdía} reports a state rather than blaming anyone; and \emph{sencillas, pero las hacemos bien} refuses the obvious conclusion.

Now read it the way somebody else would report it:

\begin{quote}
El \textbf{nivel} era \textbf{demasiado alto} para ella.
\end{quote}

Not a word of that is in what she said, and it is what she meant. She talks about the \textbf{class}; a summary talks about its \textbf{nivel}. Move between the two, because the person asking what happened will use the second.

Spanish gives you a way to leave something that is not working without criticising it, and an exam will ask you to hear the difference between \emph{it did not suit me} and \emph{it was no good}.
\end{grammarlens}

\subsection*{Your turn}
Say why she changed groups, using \emph{para mí} the way she does.

Now say it about something of your own — a class, a book, a game. Name what was too advanced, say you lost the thread, and say what you do now instead. Keep the blame out of it.

\begin{tcolorbox}[breakable,colback=teal!4,colframe=teal!35!black,title={Before you move on}]
Why did she change? (\textbf{Era demasiado avanzado para ella}.) What happened to her in class? (\textbf{Se perdía}.) Who is she with now? (\textbf{Los principiantes}.) Are the new dishes worse? (\textbf{No — sencillas, pero bien hechas}.)
\end{tcolorbox}
